\chapter*{Preface}
\addcontentsline{toc}{chapter}{Preface}

This book exists because a codebase got large enough, and honest enough about its own
gaps, that it deserved a book instead of a wiki page.

\textbf{CNA} is a C\texttt{++}23 reimplementation of the Microsoft XNA 4.0 programming
model, built on SDL3 and a pluggable graphics-backend layer. It is not a game, and not a
game engine in the ``opinionated framework with an editor'' sense --- it is a compatibility
and abstraction layer that lets code written against \texttt{Microsoft::Xna::Framework}
run as native, non-managed C\texttt{++}, across a dozen different rendering backends, on
Linux, Windows, the Web, and Android.

CNA does not exist in isolation. Around it sits a small constellation of sibling
repositories, each solving one piece of the problem so that CNA itself does not have to:
\texttt{sharp-runtime} provides a C\texttt{++} subset of the .NET base class library CNA's
API vocabulary assumes; \texttt{easy-gl} wraps OpenGL/OpenGLES behind a toolkit-independent
interface for the \texttt{EASYGL} backend; \texttt{free-direct} reimplements the narrow
slice of DirectDraw/DirectSound/DirectPlay that real legacy games actually call, backing
the \texttt{DX3} backend; \texttt{xna4-spec} is a machine-readable transcription of the
entire official XNA 4.0 API surface, used to audit CNA's coverage claims against a ground
truth rather than against memory; and \texttt{cna-samples}/\texttt{cna-examples} give the
whole thing something to actually run.

\paragraph{What this book is.} A technical reference and a guided tour, built by reading
the source code, the design documents (\texttt{NEXT.md}, \texttt{CLAUDE.md}, the \texttt{docs/}
tree, the per-backend \texttt{plan\_*.md} files), and the test suites of these repositories,
not by paraphrasing marketing claims. Where the project's own documentation is careful to
distinguish ``implemented'' from ``pixel-verified'' from ``architecturally blocked,'' this
book preserves that distinction rather than flattening it into a uniform sales pitch. CNA's
own README is unusually candid about exactly this --- it dates its own coverage numbers,
names its own known bugs, and tells you which backend's ``done'' marks mean something different
from another backend's ``done'' marks. That candor is worth preserving in a book, not smoothing over.

\paragraph{What this book is not.} It is not the official documentation of the
\texttt{openeggbert} project --- that lives in the repositories themselves, and it changes
faster than any book can. Treat this volume as a snapshot and a map: a way to build a mental
model of how the pieces fit together, which you then go verify against the current
\texttt{NEXT.md} and \texttt{docs/} tree of the actual repository, exactly as the project's
own documentation index recommends doing for itself.

\paragraph{How this volume is organized.} Part I introduces CNA and its ecosystem and the
design philosophy that runs through all of it --- API mirroring, backend/interface
separation, and an unusual commitment to differential testing against a real reference
implementation. Part II gets a reader from a clean checkout to a running triangle, and
then walks the core framework types a game is built from: the game loop, the math library,
and the content/asset model. Part III is the graphics machine in depth --- everything
between \texttt{GraphicsDevice} and a pixel actually landing on screen: \texttt{SpriteBatch},
textures and render targets, 3D models and meshes, the five stock effects, the state
objects, and the one gap CNA is most honest about, compiled \texttt{.fx} shader bytecode.
Part IV takes the same graphics machine and asks how it is actually implemented eleven
different ways, one chapter per backend, from the simplicity of \texttt{SDL\_RENDERER} to
the Wine-and-DXVK-verified authenticity of \texttt{D3D9}.

Volume II continues into input, audio, devices, networking, and gamer services; into
\texttt{sharp-runtime}, \texttt{easy-gl}, and \texttt{free-direct} as projects in their own
right; into the cross-platform engineering that makes four very different target platforms
share one test suite; and into the practical work of porting a real XNA game, using
\texttt{xna4-spec} to audit compatibility claims, and reading the project the way its own
maintainers do --- through \texttt{NEXT.md}, \texttt{plan.sqlite3}, and a Doxygen
\texttt{@note Status:} tag.

\paragraph{A note on dates.} Software documentation ages the moment it is written. Specific
percentages, task numbers, and ``as of'' dates quoted in this book (many drawn directly from
CNA's own README and \texttt{docs/} tree, dated around 2026-07-11 through 2026-07-15) are
reproduced as evidence of methodology and of the project's own verification discipline, not
as a promise that they still hold on the day you read this. Where the source material itself
flags a figure as dated or superseded, this book says so.
